\subsection{Classification Experiments}

\input{latex/figures_tables/classification_table}

Passage classification into functional roles was evaluated under two settings. The first used the original five labels in the annotation scheme, while the second merged \textit{Conclusion} into \textit{Analysis}, yielding a four-class variant. Implicit intermediate conclusions were excluded from all runs, leaving 719 explicit passages drawn from 42 cases. The four-class setting tests whether removing the sparsest label boundary simplifies the task and improves classification. In the five-class data, \textit{Conclusion} is the rarest label (44 passages), and the grouped-fold results show that merging it into \textit{Analysis} improves Macro-F1 for every model reported here.

All experiments use five-fold case-disjoint \texttt{StratifiedGroupKFold} cross-validation, grouping passages by case identifier so that no case appears in both training and test data. In the five-class experiment, the held-out folds contain 121, 162, 131, 140, and 165 passages drawn from 8, 8, 8, 9, and 9 cases, respectively. In the four-class experiment, the held-out folds contain 175, 169, 105, 129, and 141 passages drawn from 9, 9, 8, 7, and 9 cases, respectively. In both settings, the overlap between training and test cases is zero in every fold. This split design ensures that evaluation measures generalization to unseen cases rather than to additional passages from the same opinion.

We compare TF--IDF features and three embedding families: SBERT \citep{reimers2019sentencebertsentenceembeddingsusing}, LegalBERT \citep{chalkidis-etal-2020-legal}, and ModernBERT \citep{warner-etal-2025-smarter}. We also report GPT-5-mini as a separate LLM comparison row. Because its configuration differs from the embedding-based models, it should not be interpreted as a strictly like-for-like baseline.

Table~\ref{tab:classification_experiment} reports Macro-F1 and per-class F1 scores. Merging \textit{Conclusion} into \textit{Analysis} improves every model. Among the embedding-based models, LegalBERT performs best in both settings, increasing from Macro-F1 0.71 in the five-class experiment to 0.80 in the four-class experiment. GPT-5-mini is strongest in the five-class setting (0.76), but in the four-class setting LegalBERT slightly outperforms it (0.80 vs.\ 0.79). The grouped-fold baselines remain substantially lower (random: 0.17/0.22; majority: 0.13/0.18), indicating that the reported gains are not explained by label imbalance alone.

Per-label confusion analysis shows that the most important remaining boundary is \textit{Rule} versus \textit{Analysis}. In the five-class setting, GPT-5-mini predicts \textit{Rule} for 86 gold \textit{Analysis} passages (24.3\% of analyses) but predicts \textit{Analysis} for only 11 gold \textit{Rule} passages (5.3\% of rules), suggesting that it tends to overgeneralize case-specific reasoning as reusable legal standards. LegalBERT exhibits a more balanced pattern (61 \textit{Analysis}$\rightarrow$\textit{Rule} errors versus 45 \textit{Rule}$\rightarrow$\textit{Analysis} errors), whereas ModernBERT more often makes the reverse error (46 versus 61). These confusion patterns support treating \textit{Rule}/\textit{Analysis} as the central difficulty of the task and motivate reporting full confusion matrices and representative errors in an appendix.
